In this section, we introduce the basic tools needed to derive the formula
\eqref{formula}. We draw almost all the material in this section from
Section 3 of \cite{leonhardtGeometryLightScience2010}.

Recall that the main idea of cloaking is to observe the form that a
differential operator takes under a particular coordinate system and then
rewrite the equation in its original form. Differential geometry becomes
essential here, as it provides a formula for expressing the
divergence and curl operators in arbitrary coordinates. We present the formulas
in the following subsections.

\subsection{Divergence and curl}

Using differential geometry we can express the divergence in a coordinate-free
manner. For a general coordinate system with metric tensor \(g\), the
divergence operator becomes
\begin{equation}\label{divergence}
	\nabla\cdot \boldsymbol{V} = \frac{1}{\sqrt{g}}\partial_i
	\left(\sqrt{g} V^i\right).
\end{equation}

Similarly, the curl operator has the expression:
\begin{equation}\label{curl}
\left(\nabla\times \boldsymbol{V}\right)^i =\epsilon^{ijk}\partial_jV_k,
\end{equation}
where \(\epsilon^{ijk}=\pm\frac{1}{\sqrt{g}}[ijk]\) denotes the Levi-Civita tensor.

Although \eqref{divergence} and \eqref{curl} appear compact and simple,
we must interpret them carefully. First, these formulas employ a basis that has not been
normalized. Second, we must pay attention to the position of the indices:
these formulas distinguish between contravariant (upper index) vectors and
covariant (lower index) vectors. The next subsection explains this distinction in more detail.

\subsection{Raising and lowering indices}

To properly use these formulas, we must learn how to raise and lower
indices. The formulas require contravariant vectors (upper indices) in some
positions and covariant vectors (lower indices) in others. We accomplish this
using the metric tensor (to lower indices) and its inverse (to raise
indices), as follows:
\begin{align*}
	V_i&=g_{ij}V^j\\
	g^{ki}V_i&=V^k.
\end{align*}

In the first equation, we lower an index by multiplying by the metric tensor
\(g_{ij}\). In the second equation, we raise an index
by multiplying by the inverse metric tensor \(g^{ki}\).

Strictly speaking, vectors and covectors are different mathematical objects
that inhabit different spaces, but for our purposes we can treat them simply as
matrices/vectors. The above rules suffice for applying the formulas.

\subsection{Relationships between metric tensors}
The final tool we need from differential geometry is an equation that relates
the metric tensors of two different coordinate systems.

Given two coordinate systems \(\{x^i\}\) and \(\{x^{i'}\}\), we define
\[
	\Lambda=\Lambda^i_{i'}=\left[\frac{\partial x^i}{\partial
	x^{i'}}\right],
\]
the Jacobian matrix of the transformation. The relationship between the metric tensors associated with each coordinate
system is then
\begin{equation}
	\boldsymbol{g'}=\Lambda^T \boldsymbol{g} \Lambda.\label{tensorrelations}
\end{equation}
We will apply this equation in the next section.
